\documentclass[12pt, a4paper]{article}

% Pacotes essenciais
\usepackage[utf8]{inputenc}
\usepackage[T1]{fontenc}
\usepackage[portuguese]{babel}
\usepackage{amsmath, amssymb}
\usepackage{graphicx}
\usepackage{booktabs}
\usepackage{geometry}
\usepackage{hyperref}
\usepackage{float}
\usepackage{caption}
\usepackage{listings}
\usepackage{url}
\usepackage{xcolor}

\lstset{
    basicstyle=\ttfamily\footnotesize,
    keywordstyle=\color{blue},
    stringstyle=\color{red},
    commentstyle=\color{green!60!black},
    numbers=left,
    numberstyle=\tiny,
    stepnumber=1,
    numbersep=5pt,
    frame=single,
    breaklines=true,
    breakatwhitespace=false,
    showstringspaces=false,
}

\providecommand{\IBGEtab}[3]{%
    \begin{center}
        #1\vspace{0.1cm}\\
        #2\vspace{0.1cm}\\
        #3
    \end{center}
}

\providecommand{\fonte}[1]{%
    \vspace{0.1cm}
    \small\textbf{Fonte:} #1
}

\geometry{left=3cm, right=2cm, top=3cm, bottom=2cm}

\hypersetup{
    colorlinks=true,
    linkcolor=blue,
    filecolor=magenta,      
    urlcolor=blue,
    pdftitle={Pré-Projeto de Pesquisa Operacional - MECAI},
}

\title{\textbf{Projeto de Pesquisa Operacional} \\ \vspace{0.2cm} \large{ICMC USP São Carlos} \\ \vspace{0.8cm} \large Otimização da Destinação da Produção de Laranja entre Mercado Interno e Indústria de Suco Considerando Sazonalidade de Preços e Janelas Operacionais de Colheita Utilizando Programação Linear e GurobiPy}

\author{
    \textbf{Alunos:}\\
    Alexander Ferreira\\ 
    Leandro De Andrade \\ 
    \vspace{0.5cm}\\
    \textbf{Professora:} Maristela Oliveira dos Santos
}

\date{}

\begin{document}

\maketitle
\thispagestyle{empty}


\pagenumbering{arabic}

\section{Introdução}
O Brasil é o maior produtor mundial de laranja, sendo o estado de São Paulo responsável pela maior parcela da produção nacional. A região de Ribeirão Preto está inserida no principal cinturão citrícola brasileiro, apresentando forte relevância econômica para o agronegócio.

A comercialização da safra representa uma das decisões mais importantes para o produtor rural. Pequenas variações de preço entre diferentes períodos do ano podem impactar significativamente a receita obtida.

Tradicionalmente, a produção é comercializada à medida que a fruta atinge o ponto ideal de colheita. Entretanto, existe certa flexibilidade operacional no manejo do pomar que permite postergar a colheita por alguns dias, possibilitando o aproveitamento de condições comerciais mais favoráveis.

Dessa forma, técnicas de Pesquisa Operacional podem ser utilizadas para auxiliar o produtor na definição da melhor estratégia de comercialização, mapeando as alternativas de escoamento e de armazenagem temporária no pomar.

\newpage

\section{Problema de Pesquisa}
Um produtor da região de Ribeirão Preto possui uma produção prevista de 100.000 toneladas de laranja durante o segundo semestre de 2025\cite{fundecitrus2025}.

A produção pode ser destinada para:
\begin{itemize}
    \item Mercado Interno (fruta \textit{in natura});
    \item Indústria de Suco.
\end{itemize}

A laranja não pode ser armazenada após a colheita por longos períodos devido à sua elevada perecibilidade. Entretanto, a colheita pode ser postergada no pomar dentro dos seguintes limites:
\begin{itemize}
    \item Mercado Interno: até 10 dias;
    \item Indústria: até 20 dias.
\end{itemize}

O objetivo é determinar a melhor estratégia de destinação da produção para maximizar a receita da propriedade, decidindo quanto vender para o mercado, quanto para a indústria, e quanto postergar para os meses subsequentes.

\section{Objetivos}

\subsection{Objetivo Geral}
Desenvolver um modelo de Programação Linear utilizando GurobiPy para maximizar a receita líquida total de venda da produção de laranjas no segundo semestre de 2025.

\subsection{Objetivos Específicos}
\begin{itemize}
    \item Modelar matematicamente a destinação da produção;
    \item Incorporar a sazonalidade dos preços;
    \item Considerar janelas operacionais de colheita e restrições de contratos mínimos;
    \item Implementar o modelo computacionalmente em GurobiPy;
    \item Validar o modelo através da comparação com uma estratégia tradicional de venda imediata;
    \item Realizar análise de sensibilidade em diferentes cenários de mercado.
\end{itemize}

\section{Materiais e Métodos}

\subsection{Fonte dos Dados}
Os dados utilizados para a construção do modelo e balizamento da pesquisa originam-se das seguintes fontes:
\begin{itemize}
    \item \textbf{Instituto Brasileiro de Geografia e Estatística (IBGE)}: Dados referentes à produção agrícola municipal, área colhida, produtividade e valor da produção\cite{faespsenar2026laranja, fundecitrus2025}.
    \item \textbf{Fundo de Defesa da Citricultura (Fundecitrus)}: Estimativas de safra, inventário de árvores e área produtiva\cite{fundecitrus2025}.
    \item \textbf{CEPEA/ESALQ-USP}: Séries históricas e evolução sazonal de preços da laranja para a indústria e para o mercado \textit{in natura}\cite{fundecitrus2025}.
\end{itemize}

\subsection{Premissas do Estudo}
Adotou-se uma propriedade modelo com produção total considerada de 100.000 toneladas, distribuídas ao longo do semestre de acordo com a Tabela \ref{tab:producao}.

\begin{table}[H]
    \caption{Distribuição da produção no segundo semestre de 2025}
    \label{tab:producao}
    \IBGEtab{
        \begin{tabular}{lc}
        \toprule
        \textbf{Mês} & \textbf{Produção (t)} \\
        \midrule
        Julho & 25.000 \\
        Agosto & 20.000 \\
        Setembro & 18.000 \\
        Outubro & 15.000 \\
        Novembro & 12.000 \\
        Dezembro & 10.000 \\
        \midrule
        \textbf{Total} & \textbf{100.000} \\
        \bottomrule
        \end{tabular}
    }{
        \fonte{Elaborado pelos autores.}
    }{}
\end{table}

\subsection{Parâmetros de Preços Utilizados}
Os valores adotados refletem o comportamento sazonal de valorização da fruta nos períodos de entressafra e de maior demanda industrial, conforme detalhado na Tabela \ref{tab:precos}.

\begin{table}[H]
    \caption{Receita por Tonelada (R\$) - Mercado In Natura e Indústria}
    \label{tab:precos}
    \IBGEtab{
        \begin{tabular}{lcc}
        \toprule
        \textbf{Mês} & \textbf{Mercado (R\$/t)} & \textbf{Indústria (R\$/t)} \\
        \midrule
        Julho & 1.700 & 1.450 \\
        Agosto & 1.850 & 1.550 \\
        Setembro & 2.000 & 1.700 \\
        Outubro & 2.150 & 1.800 \\
        Novembro & 2.250 & 1.850 \\
        Dezembro & 2.300 & 1.900 \\
        \bottomrule
        \end{tabular}
    }{
        \fonte{Elaborado pelos autores com base em dados do CEPEA.}
    }{}
\end{table}

\section{Modelagem Matemática}

\subsection{Variáveis de Decisão}
O modelo de otimização possui quatro conjuntos principais de variáveis de decisão:
\begin{itemize}
    \item $M_t$: Quantidade vendida para o mercado interno no mês $t$.
    \item $I_t$: Quantidade vendida para a indústria de suco no mês $t$.
    \item $TM_t$: Quantidade cuja colheita é postergada do mês $t$ para o mês seguinte visando comercialização no mercado interno.
    \item $TI_t$: Quantidade cuja colheita é postergada do mês $t$ para o mês seguinte visando comercialização na indústria.
\end{itemize}

\subsection{Função Objetivo}
O objetivo do modelo é maximizar a receita total da propriedade rural, dada pela soma do volume escoado multiplicado pelo preço em cada canal e período ($T$).
\begin{equation*}
\text{Max } Z = \sum_{t \in T} PM_t M_t + \sum_{t \in T} PI_t I_t
\end{equation*}
Sendo:
\begin{itemize}
    \item $Z$: Receita total da propriedade rural.
    \item $PM_t$ e $PI_t$: Preços do mercado interno e da indústria no período $t$.
    \item $M_t$ e $I_t$: Volumes destinados a cada canal.
\end{itemize}

\subsection{Restrições}

\textbf{Balanço da Produção} \\
O escoamento e as transferências de cada mês devem ser iguais à produção daquele mês somada às transferências recebidas do mês anterior.
\begin{align*}
    &\text{Julho: } M_{\text{Jul}} + I_{\text{Jul}} + TM_{\text{Jul}} + TI_{\text{Jul}} = 25000 \\
    &\text{Agosto: } M_{\text{Ago}} + I_{\text{Ago}} + TM_{\text{Ago}} + TI_{\text{Ago}} = 20000 + TM_{\text{Jul}} + TI_{\text{Jul}} \\
    &\text{Setembro: } M_{\text{Set}} + I_{\text{Set}} + TM_{\text{Set}} + TI_{\text{Set}} = 18000 + TM_{\text{Ago}} + TI_{\text{Ago}} \\
    &\text{Outubro: } M_{\text{Out}} + I_{\text{Out}} + TM_{\text{Out}} + TI_{\text{Out}} = 15000 + TM_{\text{Set}} + TI_{\text{Set}} \\
    &\text{Novembro: } M_{\text{Nov}} + I_{\text{Nov}} + TM_{\text{Nov}} + TI_{\text{Nov}} = 12000 + TM_{\text{Out}} + TI_{\text{Out}} \\
    &\text{Dezembro: } M_{\text{Dez}} + I_{\text{Dez}} = 10000 + TM_{\text{Nov}} + TI_{\text{Nov}}
\end{align*}

\vspace{0.2cm}\noindent\textbf{Limites de Postergação no Pomar} \\
Limita o volume que pode aguardar no pé, simulando os dias máximos de tolerância da fruta.
\begin{align*}
    TM_t &\le 0.20 \cdot Q_t \quad \text{(Máximo de 10 dias para Mercado)} \\
    TI_t &\le 0.40 \cdot Q_t \quad \text{(Máximo de 20 dias para Indústria)}
\end{align*}

\vspace{0.2cm}\noindent\textbf{Restrições de Capacidade e Contrato} \\
O mercado interno absorve no máximo 12.000t/mês, e a indústria exige um contrato mínimo global de 40.000t na safra.
\begin{align*}
    M_t &\le 12000 \\
    \sum_{t \in T} I_t &\ge 40000
\end{align*}

\vspace{0.2cm}\noindent\textbf{Restrição de Não Negatividade}
\begin{equation*}
    M_t \ge 0, \quad I_t \ge 0, \quad TM_t \ge 0, \quad TI_t \ge 0
\end{equation*}

\section{Implementação Computacional e Análise de Sensibilidade}

A implementação utilizou Python, bibliotecas de manipulação de dados (Pandas, Matplotlib) e a API de otimização GurobiPy, resolvendo o problema matemático através do Algoritmo Simplex. 

\subsection{Estratégia Tradicional de Referência}
Para fins de validação do ganho de eficiência, estipulou-se a "Estratégia Tradicional", onde assume-se que toda a produção é vendida imediatamente no mês da colheita, sem postergação (transferência zero) e priorizando o canal de maior preço do mês até seu limite.
\begin{itemize}
    \item \textbf{Receita Total Tradicional:} R\$ 192.240.000,00
\end{itemize}

O indicador de ganho do modelo otimizado baseia-se na fórmula:
\begin{equation*}
    \text{Ganho (\%)} = \frac{\text{Receita Otimizada} - \text{Receita Tradicional}}{\text{Receita Tradicional}} \times 100
\end{equation*}

\subsection{Definição dos Cenários Analisados}
\begin{itemize}
    \item \textbf{Cenário 1 (Base):} Preços originais, conforme pensamento empírico do produtor. \\
    \textit{Objetivo:} Comercializar a produção sem planejamento, sem acompanhamento de mercado, tomando decisões de forma empírica.
    
    \item \textbf{Cenário 2 (Queda do Mercado Interno):} Redução de 15\% em todos os preços da fruta \textit{in natura}. \\
    \textit{Objetivo:} Avaliar a dependência da receita à comercialização in natura.
    
    \item \textbf{Cenário 3 (Aumento da Indústria):} Aumento de 20\% nos preços pagos pela indústria. \\
    \textit{Objetivo:} Avaliar a atratividade da indústria em períodos de maior demanda por suco.
    
    \item \textbf{Cenário 4 (Redução de Janelas):} Diminuição da capacidade de postergação (Mercado: de 20\% para 10\%; Indústria: de 40\% para 20\%). \\
    \textit{Objetivo:} Avaliar o impacto da redução da flexibilidade operacional da colheita sobre a receita.
\end{itemize}


\section{Resultados e Discussão}

\subsection{Desempenho por Cenário}

\textbf{Cenário 1 – Base} \\
É o cenário de referência. Como os preços do mercado interno são sempre superiores, o modelo ocupa totalmente a capacidade mensal (\textit{12.000 t/mês}), direcionando o excedente à indústria e otimizando as datas de colheita.
\begin{itemize}
    \item \textbf{Receita Máxima:} R\$ 192.240.000,00
    \item \textbf{Volumes Direcionados:} Mercado (60.000t) | Indústria (40.000t)
\end{itemize}

\vspace{0.2cm}\noindent\textbf{Cenário 2 – Queda de 15\% no Mercado Interno} \\
Com a queda, alguns meses apresentam preços do mercado interno inferiores ou muito próximos aos da indústria. O Gurobi reage reduzindo a atratividade do varejo e realocando grandes volumes para processamento industrial.
\begin{itemize}
    \item \textbf{Receita Máxima:} R\$ 173.250.000,00
    \item \textbf{Volumes Direcionados:} Mercado (48.000t) | Indústria (52.000t)
\end{itemize}

\vspace{0.2cm}\noindent\textbf{Cenário 3 – Aumento de 20\% na Indústria} \\
A indústria torna-se altamente rentável. O modelo rompe os padrões de abastecimento primário do mercado e concentra pesadamente as vendas para as fábricas de suco, indicando uma maior demanda do produto pela varejo.
\begin{itemize}
    \item \textbf{Receita Máxima:} R\$ 206.088.000,00
    \item \textbf{Volumes Direcionados:} Mercado (24.000t) | Indústria (76.000t)
\end{itemize}

\vspace{0.2cm}\noindent\textbf{Cenário 4 – Redução das Janelas Operacionais} \\
O produtor perde flexibilidade para reter a laranja no pé reduzindo a janela operacional em 50\% (Varejo passa de 10 dias para 5 dias e Industria de 20 dias para 10 dias). Consequentemente, não consegue surfar os picos de preço de final de ano na mesma proporção.
\begin{itemize}
    \item \textbf{Receita Máxima:} R\$ 189.645.000,00
    \item \textbf{Volumes Direcionados:} Mercado (60.000t) | Indústria (40.000t)
\end{itemize}

\subsection{Comparativo Geral Otimizado}

A consolidação de todos os cenários pode ser visualizada na Tabela \ref{tab:comparativo} e na Figura \ref{fig:receita_maxima}.

\begin{table}[H]
    \caption{Comparativo Consolidado de Receitas e Volumes por Cenário}
    \label{tab:comparativo}
    \IBGEtab{
        \resizebox{\textwidth}{!}{%
            \begin{tabular}{lcccc}
            \toprule
            \textbf{Cenário} & \textbf{Receita Total (R\$)} & \textbf{Mercado (t)} & \textbf{Indústria (t)} & \textbf{$\Delta$ Receita (\%)} \\
            \midrule
            1. Base & 192.240.000,00 & 60.000 & 40.000 & 0,00 \\
            2. Mercado -15\% & 173.250.000,00 & 48.000 & 52.000 & -9,87 \\
            3. Indústria +20\% & 206.088.000,00 & 24.000 & 76.000 & +7,20 \\
            4. Janela Reduzida & 189.645.000,00 & 60.000 & 40.000 & -1,34 \\
            \bottomrule
            \end{tabular}%
        }
    }{
        \fonte{Elaborado pelos autores - Comparativo dos Cenários.}
    }{}
\end{table}

\begin{figure}[H]
    \centering
    \caption{Receita Máxima por Cenário (em milhões de R\$)}
    \vspace{0.2cm}
    \includegraphics[width=0.75\linewidth]{grafico.png} \\
    \fonte{Elaborado pelos autores - Comparativo dos Cenários.}
    \label{fig:receita_maxima}
\end{figure}

\subsection{Principais Indicadores Obtidos}
\begin{itemize}
    \item \textbf{Cenário Base:} Base (R\$ 192.240.000,00).
    \item \textbf{Cenário Mercado -15\%:} Mercado -15\% (R\$ 173.250.000,00).
    \item \textbf{Cenário - melhor identificado:} Indústria +20\% (R\$ 206.088.000,00).
    \item \textbf{Cenário Janela Reduzida:} Janela Reduzida\% (R\$ 189.645.000,00).\\[0.3cm]
    \textbf{Ganho de Otimização:} O melhor cenário alcançado pela Programação Linear reflete um ganho puro de \textbf{+7,20\%} de rentabilidade global em relação à Estratégia Tradicional de venda empírico do produtor.
\end{itemize}

\section{Conclusões}
O conjunto dos cenários evidenciou que a receita citrícola é sensível às variações de mercado e às restrições operacionais, e o modelo de Programação Linear proposto, resolvido pelo algoritmo Simplex, mostrou-se capaz de representar de forma realista o processo decisório, com a perecibilidade da fruta e as janelas de postergação no pomar incorporadas de forma explícita, oferecendo suporte quantitativo à estratégia ótima de comercialização. \\[0.05cm]

As simulações evidenciam dinâmicas cruciais:
\begin{itemize}
    \item A flexibilização da colheita tem um valor econômico mensurável. No Cenário 4 (Janela Reduzida) a receita cai para R\$ 189,65 milhões, uma perda R\$ 2,59 (-1,34\%) em relação a Base. Isso mostra que reduzir as janelas de postergação limita a capacidade de capturar preços melhores. 
     \item O mercado interno é vulnerável a quedas de preços. No Cenário 2 (Cenário Mercado -15\%) a receita despenca para R\$ 173,25 milhões, uma redução de R\$ 19 milhões (-9,87\%) em relação a Base.
    \item A rentabilidade é altamente sensível ao preço pago. No Cenário 3 (Indústria +20\%) a receita sobe para R\$ 206,09 milhões, um ganho de R\$ 13,85 milhões (+7,20\%) em relação a Base. O volume destinado à indústria aumenta para 76.000t, evidenciando que esse canal se torna prioritário quando valorizado. Com isso, o modelo otimizado supera a estratégia tradicional, empírica do produtor.  
\end{itemize}

A aplicação prática do algoritmo Simplex via GurobiPy demonstrou-se robusta, fornecendo uma ferramenta gerencial objetiva. Os resultados provam que a Pesquisa Operacional aplicada ao agronegócio permite transformar a intuição empírica da comercialização de commodities em um modelo quantitativo altamente eficiente, garantindo a maximização do faturamento rural.\\[0.03cm]

Para garantir transparência e permitir a consulta e a reprodução do trabalho, o projeto foi disponibilizado publicamente no GitHub, podendo ser acessado pelo seguinte endereço:


\begin{center}
\href{https://github.com/alexferrei198170/Projeto-Pesquisa-Operacional-Otimiza-o-Destina-o-Produ-o-Laranja-}{\textbf{Repositório do projeto no GitHub}}
\end{center}

% Este é um comentário PAREI=FALTA ALTERAR O CODIGO E REFERENCIAS

\newpage
\appendix
\section{Código Fonte - GurobiPy}
\begin{lstlisting}[language=Python, title=Otimização da Destinação da Produção de Laranja entre Mercado Interno e Indústria]
from gurobipy import *
import pandas as pd
import matplotlib.pyplot as plt

# FUNÇÃO DE OTIMIZAÇÃO
def executar_modelo(
    nome_cenario,
    preco_mercado,
    preco_industria,
    perc_mercado=0.20,
    perc_industria=0.40
):
    modelo = Model(nome_cenario)
    modelo.setParam("OutputFlag", 0)
    meses = [
        "Julho",
        "Agosto",
        "Setembro",
        "Outubro",
        "Novembro",
        "Dezembro"
    ]
    producao = {
        "Julho": 25000,
        "Agosto": 20000,
        "Setembro": 18000,
        "Outubro": 15000,
        "Novembro": 12000,
        "Dezembro": 10000
    }
    # VARIÁVEIS DE DECISÃO
    M = modelo.addVars(meses, lb=0, name="Mercado")
    I = modelo.addVars(meses, lb=0, name="Industria")
    TM = modelo.addVars(
        meses[:-1],
        lb=0,
        name="TransferenciaMercado"
    )
    TI = modelo.addVars(
        meses[:-1],
        lb=0,
        name="TransferenciaIndustria"
    )
    # FUNÇÃO OBJETIVO
    modelo.setObjective(
        quicksum(
            preco_mercado[m] * M[m]
            for m in meses
        )
        +
        quicksum(
            preco_industria[m] * I[m]
            for m in meses
        ),
        GRB.MAXIMIZE
    )
    # BALANÇO DE PRODUÇÃO
    # Julho
    modelo.addConstr(
        M["Julho"]
        +
        I["Julho"]
        +
        TM["Julho"]
        +
        TI["Julho"]
        ==
        producao["Julho"]
    )
    # Agosto até Novembro
    for i in range(1, 5):
        atual = meses[i]
        anterior = meses[i - 1]
        modelo.addConstr(
            M[atual]
            +
            I[atual]
            +
            TM[atual]
            +
            TI[atual]
            ==
            producao[atual]
            +
            TM[anterior]
            +
            TI[anterior]
        )
    # Dezembro
    modelo.addConstr(
        M["Dezembro"]
        +
        I["Dezembro"]
        ==
        producao["Dezembro"]
        +
        TM["Novembro"]
        +
        TI["Novembro"]
    )

    # POSTERGAÇÃO DA COLHEITA
    for mes in meses[:-1]:
        modelo.addConstr(
            TM[mes]
            <=
            perc_mercado * producao[mes]
        )
        modelo.addConstr(
            TI[mes]
            <=
            perc_industria * producao[mes]
        )
    # CAPACIDADE DO MERCADO INTERNO
    for mes in meses:
        modelo.addConstr(
            M[mes]
            <=
            12000
        )
    # CONTRATO MÍNIMO COM A INDÚSTRIA
    modelo.addConstr(
        quicksum(
            I[m]
            for m in meses
        )
        >=
        40000
    )
    # RESOLUÇÃO
    modelo.optimize()
    if modelo.Status != GRB.OPTIMAL:
        raise Exception(
            f"Não foi encontrada solução ótima para {nome_cenario}"
        )
    receita = modelo.ObjVal
    mercado_total = sum(
        M[m].X
        for m in meses
    )
    industria_total = sum(
        I[m].X
        for m in meses
    )
    return {
        "receita": receita,
        "mercado": mercado_total,
        "industria": industria_total
    }
# RECEITA TRADICIONAL (SEM POSTERGAÇÃO)
receita_tradicional = (
    (12000 * 1700 + 13000 * 1450)
    +
    (12000 * 1850 + 8000 * 1550)
    +
    (12000 * 2000 + 6000 * 1700)
    +
    (12000 * 2150 + 3000 * 1800)
    +
    (12000 * 2250)
    +
    (10000 * 2300)
)
# PREÇOS 
mercado_base = {
    "Julho": 1700,
    "Agosto": 1850,
    "Setembro": 2000,
    "Outubro": 2150,
    "Novembro": 2250,
    "Dezembro": 2300
}
industria_base = {
    "Julho": 1450,
    "Agosto": 1550,
    "Setembro": 1700,
    "Outubro": 1800,
    "Novembro": 1850,
    "Dezembro": 1900
}
# CENÁRIO 1 # BASE
cenario1 = executar_modelo(
    "Base",
    mercado_base,
    industria_base
)
# CENÁRIO 2 # MERCADO -15%
mercado_c2 = {
    k: v * 0.85
    for k, v in mercado_base.items()
}
cenario2 = executar_modelo(
    "Mercado -15%",
    mercado_c2,
    industria_base
)
# CENÁRIO 3 # INDÚSTRIA +20%
industria_c3 = {
    k: v * 1.20
    for k, v in industria_base.items()
}
cenario3 = executar_modelo(
    "Indústria +20%",
    mercado_base,
    industria_c3
)
# CENÁRIO 4 # REDUÇÃO DAS JANELAS
cenario4 = executar_modelo(
    "Janela Reduzida",
    mercado_base,
    industria_base,
    perc_mercado=0.10,
    perc_industria=0.20
)
# COMPARATIVO
comparativo = pd.DataFrame({
    "Cenário": [
        "Base",
        "Mercado -15%",
        "Indústria +20%",
        "Janela Reduzida"
    ],
    "Receita": [
        cenario1["receita"],
        cenario2["receita"],
        cenario3["receita"],
        cenario4["receita"]
    ],
    "Mercado Interno (t)": [
        cenario1["mercado"],
        cenario2["mercado"],
        cenario3["mercado"],
        cenario4["mercado"]
    ],
    "Indústria (t)": [
        cenario1["industria"],
        cenario2["industria"],
        cenario3["industria"],
        cenario4["industria"]
    ]
})
comparativo["Variação (%)"] = (
    (
        comparativo["Receita"]
        -
        cenario1["receita"]
    )
    /
    cenario1["receita"]
) * 100
print("\n")
print("=" * 100)
print("COMPARATIVO DOS CENÁRIOS")
print("=" * 100)
print(comparativo)
# GRÁFICO
plt.figure(figsize=(10,6))
# Cores específicas para cada cenário
cores = ["orange","gold","green","blue"]
# Barras
plt.bar(comparativo["Cenário"], comparativo["Receita"]/1e6, color=cores)
# Título e eixos
plt.title("Receita Máxima por Cenário")
plt.ylabel("Receita (Milhões de R$)")
plt.grid(axis="y", alpha=0.3)
# Valores acima das barras
for i, valor in enumerate(comparativo["Receita"]):
    plt.text(i, valor/1e6 + 0.5, f"{valor/1e6:.2f}", ha="center")
plt.tight_layout()
plt.show()
# CONCLUSÕES AUTOMÁTICAS
melhor = comparativo.loc[
    comparativo["Receita"].idxmax()
]
print("\n")
print("=" * 100)
print("CONCLUSÕES")
print("=" * 100)
print(
    f"Melhor Cenário: "
    f"{melhor['Cenário']}"
)
print(
    f"Receita Máxima: "
    f"R$ {melhor['Receita']:,.2f}"
)
print("\nINTERPRETAÇÃO")
print(
    "- O modelo identifica a alocação ótima entre mercado interno e indústria."
)
print(
    "- A postergação limitada da colheita possui valor econômico."
)
print(
    "- A indústria torna-se mais atrativa quando seus preços aumentam."
)
print(
    "- A redução das janelas operacionais diminui a flexibilidade do sistema."
)
print(
    "- A Programação Linear mostrou-se eficaz como ferramenta de apoio à decisão."
)
\end{lstlisting}

\newpage
\begin{thebibliography}{99}

\bibitem{boldrini1984}
J. L. Boldrini et al. \textit{Álgebra Linear e Aplicações}. Harbra, 1984.

\bibitem{cuminato2014}
J. A. Cuminato and V. Ruas. Unification of distance inequalities. \textit{Computational and Applied Mathematics}, 2014.

\bibitem{jafelice2014}
R. M. Jafelice et al. Study of the dynamics of HIV. \textit{Computational and Applied Mathematics}, 2014.

\bibitem{santos2013}
I. L. D. Santos e G. N. Silva. Controle ótimo em escalas temporais. \textit{SBMAC}, 2013.

\bibitem{cepea2026}
CEPEA – Centro de Estudos Avançados em Economia Aplicada. Indicadores agropecuários: Citros, 2026.

\bibitem{fundecitrus2025}
Fundecitrus. Inventário de árvores de citros do cinturão citrícola de São Paulo e triângulo/sudoeste mineiro. \textit{Technical report}, Fundo de Defesa da Citricultura, Araraquara, SP, 2025. Acessado em: 16 set. 2026.

\end{thebibliography}

%================================================================================
% Referências Bibliográficas
%\bibliographystyle{plain}
%\bibliography{referencias}
% Arquivo contendo as referências
%================================================================================

\end{document}




% Referências Bibliográficas

\begin{thebibliography}{99}
\bibitem{cepea} CEPEA - Centro de Estudos Avançados em Economia Aplicada. Indicadores agropecuários: Citros, 2026.
\bibitem{faesp} Federação da Agricultura e do Serviço Nacional de Aprendizagem Rural de São Paulo - FAESP/SENAR. Laranja paulista aproxima culturas e conecta o mundo, 2026.
\bibitem{ferreira2026} FERREIRA, Alexander; NESPOLO, Renan Guilherme. Previsão de commodity laranja utilizando métodos de regressão e ciência de dados. \textit{Interface Tecnológica}, 22(2):303-315, 2026.
\bibitem{fundecitrus2025} FUNDECITRUS. Inventário de árvores de citros do cinturão citrícola de São Paulo e Triângulo/Sudoeste Mineiro. Relatório Técnico, Fundo de Defesa da Citricultura, Araraquara, SP, 2025.
\end{thebibliography}